\chapter{Many-Electron Atoms and Term Symbols}\label{ch:b3:many-electron-atoms}

In 1868 a yellow line in the spectrum of the Sun revealed an element
unknown on Earth, helium. When it was finally isolated in the laboratory, in
1895, its spectrum looked like that of \emph{two} elements: two families of
lines, which for a while were attributed to ``orthohelium'' and
``parahelium'', and which never exchanged light with each other. There is
only one helium. The two families are the triplet and \omterm{def:b3:many-electron-atoms:singlet-triplet}{singlet states} of the
same atom, kept apart by the electron's spin and by a rule more fundamental
than any force: the wavefunction of two electrons must change sign when they
are exchanged. This chapter follows that rule from the Pauli principle to
the \omterm{def:b3:many-electron-atoms:term}{term symbols} with which spectroscopists label every state of every atom
and ion --- including the transition-metal ions whose colours are explained
in \cref{ch:b3:complex-spectra-magnetism}.

\begin{recall}
The Year 1 volume labelled electrons by four quantum numbers $n$, $l$,
$m_l$, $m_s$, built ground configurations with the Pauli principle, the
Klechkowski rule and Hund's rule, and counted unpaired electrons. The Year 2
volume gave orbital energies and Slater's rules for screening.
\Cref{ch:b3:quantum-model-systems} supplied \omterm{def:b3:quantum-model-systems:operator}{operators}, \omterm{def:b3:quantum-model-systems:operator}{eigenvalues}, the
rotor's angular momentum ($\hat L^2$ with \omterm{def:b3:quantum-model-systems:operator}{eigenvalues} $l(l+1)\hbar^2$) and the
hydrogen atom.
\end{recall}

\section{Spin and spin-orbitals}

The electron carries an intrinsic angular momentum, its spin, with quantum
number $s = \frac12$: $\hat S^2$ has the single \omterm{def:b3:quantum-model-systems:operator}{eigenvalue} $s(s+1)\hbar^2 =
\frac34\hbar^2$, and $\hat S_z$ the two \omterm{def:b3:quantum-model-systems:operator}{eigenvalues} $m_s\hbar = \pm\frac12\hbar$.
The two spin functions are written $\alpha$ ($m_s = +\frac12$) and $\beta$
($m_s = -\frac12$); they are orthonormal, $\langle\alpha|\alpha\rangle =
\langle\beta|\beta\rangle = 1$ and $\langle\alpha|\beta\rangle = 0$. Spin has no
classical counterpart; it is not the rotation of a small sphere, and it
does not enter a non-relativistic Hamiltonian at all.

\begin{definition}[Spin-orbital]\label{def:b3:many-electron-atoms:spin-orbital}
A \emph{spin-orbital}\index{spin-orbital} is the product of a spatial orbital
$\phi(\mathbf r)$ and a spin function: $\phi\alpha$ or $\phi\beta$. A
spin-orbital holds the complete description of one electron.
\end{definition}

An atomic orbital labelled $(n, l, m_l)$ in the Year 1 volume thus gives two
\omterm{def:b3:many-electron-atoms:spin-orbital}{spin-orbitals}, which is why a subshell holds $2(2l+1)$ electrons. The real
question is why a \omterm{def:b3:many-electron-atoms:spin-orbital}{spin-orbital} holds at most one.

\section{Indistinguishable electrons and the Slater determinant}

\begin{definition}[Indistinguishable particles, antisymmetric wavefunction]\label{def:b3:many-electron-atoms:indistinguishable}
Particles are \emph{indistinguishable}\index{indistinguishable particles} if
no measurement can tell them apart: every electron is identical to every
other. A wavefunction of several electrons is an \emph{antisymmetric
wavefunction}\index{antisymmetric wavefunction} if it changes sign when the
coordinates (space and spin) of any two electrons are exchanged:
$\Psi(\dots,i,\dots,j,\dots) = -\Psi(\dots,j,\dots,i,\dots)$.
\end{definition}

Indistinguishability alone requires $|\Psi|^2$ to be unchanged by an
exchange, so $\Psi$ can at most change sign. Nature has chosen: the
wavefunction of any system of electrons --- of any fermions --- is
antisymmetric. This is a postulate of quantum mechanics, confirmed by every
atomic spectrum.

\begin{definition}[Slater determinant]\label{def:b3:many-electron-atoms:slater}
For $N$ electrons in $N$ different \omterm{def:b3:many-electron-atoms:spin-orbital}{spin-orbitals} $\chi_1, \dots, \chi_N$, the
\emph{Slater determinant}\index{Slater determinant} is
\[
\Psi = \frac{1}{\sqrt{N!}}\begin{vmatrix}
\chi_1(1) & \chi_2(1) & \cdots & \chi_N(1)\\
\chi_1(2) & \chi_2(2) & \cdots & \chi_N(2)\\
\vdots & & & \vdots\\
\chi_1(N) & \chi_2(N) & \cdots & \chi_N(N)
\end{vmatrix},
\]
where row $i$ holds electron $i$ in each \omterm{def:b3:many-electron-atoms:spin-orbital}{spin-orbital}. It is written for short
$|\chi_1\chi_2\dots\chi_N|$.
\end{definition}

\begin{theorem}[The Pauli principle]\label{thm:b3:many-electron-atoms:pauli}
A \omterm{def:b3:many-electron-atoms:slater}{Slater determinant} is antisymmetric, and it vanishes if two electrons
occupy the same \omterm{def:b3:many-electron-atoms:spin-orbital}{spin-orbital}. Hence no two electrons of an atom can have the
same four quantum numbers.
\end{theorem}

\begin{proof}
Exchanging electrons $i$ and $j$ exchanges two rows of the determinant,
which changes its sign: antisymmetry holds for every pair. If two
\omterm{def:b3:many-electron-atoms:spin-orbital}{spin-orbitals} are equal, two columns are equal and the determinant is zero:
no such state exists. Two electrons of the same atom with the same $n$, $l$,
$m_l$ and $m_s$ would sit in the same \omterm{def:b3:many-electron-atoms:spin-orbital}{spin-orbital}.
\end{proof}

\begin{example}[The ground state of helium]\label{ex:b3:many-electron-atoms:helium-ground}
With both electrons in $1s$,
\[
\Psi = \frac{1}{\sqrt2}\begin{vmatrix}1s\alpha(1) & 1s\beta(1)\\ 1s\alpha(2) &
1s\beta(2)\end{vmatrix} = 1s(1)\,1s(2)\cdot\frac{1}{\sqrt2}\big[\alpha(1)\beta(2)
- \beta(1)\alpha(2)\big].
\]
The spatial part is symmetric, the spin part antisymmetric: the two spins
are paired, with total spin zero.
\end{example}

\section{Helium: Coulomb repulsion and exchange}

Excite one electron of helium to $2s$. Four determinants can be built from
$1s\alpha$, $1s\beta$, $2s\alpha$, $2s\beta$ with one electron in each orbital.
Rearranged, they become products of one spatial and one spin function:
\[
\Psi_\pm = \frac{1}{\sqrt2}\big[1s(1)2s(2) \pm 2s(1)1s(2)\big]\times(\text{spin
function}).
\]
The antisymmetric spatial function $\Psi_-$ must be paired with a symmetric
spin function, of which there are three: $\alpha(1)\alpha(2)$,
$\beta(1)\beta(2)$ and $\frac{1}{\sqrt2}[\alpha(1)\beta(2) + \beta(1)\alpha(2)]$, the
three components $M_S = 1, 0, -1$ of a total spin $S = 1$. The symmetric
$\Psi_+$ must be paired with the single antisymmetric spin function, $S = 0$.

\begin{definition}[Singlet and triplet states]\label{def:b3:many-electron-atoms:singlet-triplet}
A state of total spin $S = 0$ is a \emph{singlet state}\index{singlet state};
one of total spin $S = 1$, whose three components $M_S = 1, 0, -1$ have the
same energy in the absence of a magnetic field, is a \emph{triplet
state}\index{triplet state}.
\end{definition}

\begin{definition}[Exchange integral]\label{def:b3:many-electron-atoms:exchange}
For two orbitals $a$ and $b$ and the electron repulsion $e^2/4\pi\varepsilon_0
r_{12}$, the \emph{exchange integral}\index{exchange integral} is
\[
K_{ab} = \iint a(1)b(2)\,\frac{e^2}{4\pi\varepsilon_0r_{12}}\,b(1)a(2)\,\dd\tau_1\dd\tau_2,
\]
the same integral as the classical electron repulsion $J_{ab}$ (with
$a(1)b(2)$ on both sides) except that the electrons are swapped on one
side. $K_{ab}$ is positive and has no classical counterpart.
\end{definition}

\begin{proposition}[Singlet and triplet energies]\label{prop:b3:many-electron-atoms:helium}
To first order in the electron repulsion, the $1s2s$ states of helium have
energies $E_\pm = E_{1s} + E_{2s} + J \pm K$, the singlet ($+$) above the
triplet ($-$) by $2K$.
\end{proposition}

\begin{proof}
The unperturbed energy of both $\Psi_\pm$ is $E_{1s} + E_{2s}$ (hydrogen-like
orbitals of nuclear charge 2). The first-order correction is the
\omterm{def:b3:quantum-model-systems:hermitian}{expectation value} of the repulsion $\hat V_{12}$ in the normalised state:
$\frac12\langle 1s(1)2s(2) \pm 2s(1)1s(2)|\hat V_{12}|1s(1)2s(2) \pm
2s(1)1s(2)\rangle$. The two direct terms give $\frac12(J + J)$; the two cross
terms give $\pm\frac12(K + K)$, since $\hat V_{12}$ is symmetric in the
electrons. So $\Delta E = J \pm K$, and the spin functions, normalised and
untouched by $\hat V_{12}$, only multiply by 1.
\end{proof}

The triplet lies lower although the Hamiltonian contains no spin: its
spatial function vanishes when $\mathbf r_1 = \mathbf r_2$, so the two
electrons avoid each other and repel less. This ``Fermi hole'' is the
physical content of Hund's first rule.

\begin{example}[The exchange integral of helium, measured]\label{ex:b3:many-electron-atoms:K}
% ledger: lev:He.2s3S1, lev:He.2s1S0
The $1s2s$ levels of helium lie at \qty{159856}{cm^{-1}} (${}^3$S) and
\qty{166277}{cm^{-1}} (${}^1$S) above the ground state: the singlet is
\qty{6421}{cm^{-1}}, \qty{0.796}{eV}, higher, and $K(1s,2s) = \qty{0.398}{eV}$.
The figure below shows the two families.
\end{example}

\section{Russell--Saunders coupling and term symbols}

In a light atom, the orbital angular momenta of the electrons add up to a
total $\mathbf L$, their spins to a total $\mathbf S$, and only then do
$\mathbf L$ and $\mathbf S$ couple weakly with each other.

\begin{definition}[Russell--Saunders coupling]\label{def:b3:many-electron-atoms:russell-saunders}
In \emph{Russell--Saunders coupling}\index{Russell--Saunders coupling}, the
states of a configuration are labelled by the \emph{total orbital angular
momentum quantum number}\index{total orbital angular momentum quantum number}
$L$ ($\hat{\mathbf L}^2 = L(L+1)\hbar^2$, with $M_L = \sum m_l$), the \emph{total
spin quantum number}\index{total spin quantum number} $S$ ($M_S = \sum m_s$),
and the \emph{total angular momentum quantum number}\index{total angular momentum quantum number}
$J$, which takes the values $|L - S|, \dots, L + S$.
\end{definition}

\begin{definition}[Spectroscopic term, spin multiplicity, term symbol]\label{def:b3:many-electron-atoms:term}
A \emph{spectroscopic term}\index{spectroscopic term} is the set of the
$(2L+1)(2S+1)$ states of a configuration with given $L$ and $S$. Its
\emph{spin multiplicity}\index{spin multiplicity} is $2S + 1$. The \emph{term symbol}
\index{term symbol} is ${}^{2S+1}L$, with the letters S, P, D, F, G, H
for $L = 0, 1, 2, 3, 4, 5$; a state of given $J$ is written ${}^{2S+1}L_J$ ---
for example \termsym{3}{P}{2}.
\end{definition}

\begin{proposition}[Closed shells]\label{prop:b3:many-electron-atoms:closed-shell}
A filled subshell has $L = 0$ and $S = 0$; it contributes nothing to the
\omterm{def:b3:many-electron-atoms:term}{term symbol}.
\end{proposition}

\begin{proof}
In a filled subshell every $m_l$ appears twice and every $m_s = \pm\frac12$
appears $2l+1$ times, so $M_L = 0$ and $M_S = 0$; and there is only one way to
fill it (one determinant). A single state with $M_L = M_S = 0$ can only
belong to $L = 0$, $S = 0$.
\end{proof}

\begin{proposition}[Counting the states of a configuration]\label{prop:b3:many-electron-atoms:count}
$N$ electrons in a subshell of $2(2l+1)$ \omterm{def:b3:many-electron-atoms:spin-orbital}{spin-orbitals} give
$\binom{2(2l+1)}{N}$ determinants, and the degeneracies $(2L+1)(2S+1)$ of
its terms add up to this number.
\end{proposition}

\begin{proof}
A determinant is fixed by the set of occupied \omterm{def:b3:many-electron-atoms:spin-orbital}{spin-orbitals} (their order
only changes its sign): one chooses $N$ of the $2(2l+1)$. The terms are the
same states regrouped, so their counts must agree.
\end{proof}

\begin{method}[Terms of a configuration]\label{met:b3:many-electron-atoms:terms}
\begin{enumerate}
  \item List the determinants allowed by the Pauli principle and tabulate
        them by $M_L$ and $M_S$.
  \item Find the largest $M_L$; with the largest $M_S$ that accompanies it,
        it starts a term with $L = M_L$, $S = M_S$.
  \item Strike out one determinant for each pair $(M_L, M_S)$ with
        $|M_L| \le L$, $|M_S| \le S$.
  \item Repeat with what is left, until the table is empty; check the count.
\end{enumerate}
\end{method}

\begin{example}[The terms of $p^2$]\label{ex:b3:many-electron-atoms:p2}
Two electrons in $p$ ($m_l = 1, 0, -1$) give $\binom62 = 15$ determinants.
Their table by $M_L$ and $M_S$ is:
\begin{center}\small
\begin{tabular}{c|ccc}
\hline
 & $M_S = 1$ & $M_S = 0$ & $M_S = -1$\\ \hline
$M_L = 2$ & -- & $(1^+,1^-)$ & --\\
$M_L = 1$ & $(1^+,0^+)$ & $(1^+,0^-)$, $(1^-,0^+)$ & $(1^-,0^-)$\\
$M_L = 0$ & $(1^+,-1^+)$ & $(1^+,-1^-)$, $(1^-,-1^+)$, $(0^+,0^-)$ & $(1^-,-1^-)$\\
$M_L = -1$ & $(0^+,-1^+)$ & $(0^+,-1^-)$, $(0^-,-1^+)$ & $(0^-,-1^-)$\\
$M_L = -2$ & -- & $(-1^+,-1^-)$ & --\\ \hline
\end{tabular}
\end{center}
(with $m_l^\pm$ for $m_s = \pm\frac12$). $M_L = 2$ occurs only with $M_S = 0$:
a ${}^1$D term (5 states). The largest remaining $M_L = 1$ comes with
$M_S = 1$: a ${}^3$P term (9 states). One state with $M_L = M_S = 0$ remains: a
${}^1$S term. Indeed $5 + 9 + 1 = 15$.
\end{example}

\begin{example}[Carbon, measured]\label{ex:b3:many-electron-atoms:carbon}
% ledger: lev:C.3P1, lev:C.3P2, lev:C.1D2, lev:C.1S0
The ground configuration $2p^2$ of carbon gives the levels \termsym{3}{P}{0},
\termsym{3}{P}{1}, \termsym{3}{P}{2} at 0, 16.4 and \qty{43.4}{cm^{-1}}, then
\termsym{1}{D}{2} at \qty{10193}{cm^{-1}} and \termsym{1}{S}{0} at
\qty{21648}{cm^{-1}}: the ${}^3$P term is lowest, as Hund's rules predict.
\end{example}

\begin{omfigure}
\begin{tikzpicture}
\begin{axis}[name=main, width=0.62\linewidth, height=6.4cm, xmin=-0.3, xmax=6.6,
  ymin=-1500, ymax=23500, axis y line=left, axis x line=none,
  ylabel={energy / cm$^{-1}$}, unbounded coords=jump, clip=false,
  scaled y ticks=false, ytick={0,10000,20000},
  tick label style={font=\small}, label style={font=\small}]
\addplot[omstate] table[x=x, y=E] {figdata/out/bachelor-3-atomic-levels-carbon.dat};
\draw[densely dotted, gray] (axis cs:1,4858.5) -- (axis cs:2,29.6);
\draw[densely dotted, gray] (axis cs:1,4858.5) -- (axis cs:2,10192.7);
\draw[densely dotted, gray] (axis cs:1,4858.5) -- (axis cs:2,21648);
\draw[densely dotted, gray] (axis cs:3,10192.7) -- (axis cs:4,10192.7);
\draw[densely dotted, gray] (axis cs:3,21648) -- (axis cs:4,21648);
\node[font=\small, anchor=south] at (axis cs:0.5,4858.5) {$2p^2$};
\node[font=\small, anchor=south] at (axis cs:2.5,21648) {${}^1$S};
\node[font=\small, anchor=south] at (axis cs:2.5,10192.7) {${}^1$D};
\node[font=\small, anchor=south] at (axis cs:2.5,29.6) {${}^3$P};
\node[font=\small, anchor=west] at (axis cs:5.05,21648) {\termsym{1}{S}{0}};
\node[font=\small, anchor=west] at (axis cs:5.05,10192.7) {\termsym{1}{D}{2}};
\node[font=\scriptsize, anchor=north] at (axis cs:0.5,-600) {configuration};
\node[font=\scriptsize, anchor=north] at (axis cs:2.5,-600) {terms};
\node[font=\scriptsize, anchor=north] at (axis cs:4.5,-600) {levels};
\draw[gray, ->] (axis cs:3.05,60) -- (axis cs:6.9,2500);
\end{axis}
\begin{axis}[at={(main.south east)}, anchor=south west, xshift=1.2cm, yshift=0.9cm,
  width=3.6cm, height=3.6cm, xmin=-0.1, xmax=1.1, ymin=-5, ymax=50,
  axis y line=right, axis x line=none, unbounded coords=jump, ytick={0,16.4,43.4},
  yticklabels={0,16.4,43.4}, tick label style={font=\scriptsize}, clip=false,
  title={\scriptsize ${}^3$P, zoomed}, title style={yshift=-4pt}]
\addplot[omstate] table[x=x, y=E] {figdata/out/bachelor-3-atomic-levels-carbon-zoom.dat};
\node[font=\scriptsize, anchor=east] at (axis cs:0,0) {$J = 0$};
\node[font=\scriptsize, anchor=east] at (axis cs:0,16.4) {$J = 1$};
\node[font=\scriptsize, anchor=east] at (axis cs:0,43.4) {$J = 2$};
\end{axis}
\end{tikzpicture}

{\small The ground configuration of carbon: one configuration, three terms
(at their measured energies; the ${}^3$P term at the weighted mean of its
levels, the configuration at the mean of all fifteen states) and five
levels. The spin--orbit splitting of ${}^3$P (zoom, in
\unit{cm^{-1}}) is a few hundred times smaller than the separations between
terms.}
\end{omfigure}

\section{Hund's rules and spin--orbit coupling}

\begin{proposition}[Hund's rules for terms]\label{prop:b3:many-electron-atoms:hund-terms}
For the ground configuration of an atom or ion, the lowest term is the one
with (1) the largest $S$, then (2) for that $S$, the largest $L$. Its lowest
level has (3) $J = |L - S|$ if the subshell is less than half full, $J = L +
S$ if it is more than half full.
\end{proposition}
\begin{proof}[Status]
These rules are established experimentally, for ground configurations only;
they are not theorems. Rule 1 is the exchange stabilisation of the previous
section; rule 2 says that electrons turning in the same direction meet less
often; rule 3 follows from the sign of the \omterm{def:b3:many-electron-atoms:spin-orbit}{spin--orbit coupling}, below.
\end{proof}

\begin{method}[The ground term from a box diagram]\label{met:b3:many-electron-atoms:ground-term}
\begin{enumerate}
  \item Fill the boxes of the open subshell from $m_l = +l$ downwards, one
        electron per box with spins parallel, then pair from $+l$ again.
  \item $S = \frac12 \times$(number of unpaired electrons); $L = |\sum m_l|$.
  \item $J = |L - S|$ for less than half full, $L + S$ for more than half
        full; for exactly half full $L = 0$ and $J = S$.
\end{enumerate}
\end{method}

\begin{example}[Ground terms]\label{ex:b3:many-electron-atoms:ground-terms}
% ledger: lev:Ti2.3P0, lev:Ni2.3P0, lev:Cr3.4P1_2
\ce{Ti^2+}, $d^2$: boxes $m_l = 2, 1$ filled, $S = 1$, $L = 3$, $J = 2$:
\termsym{3}{F}{2}. \ce{Cr^3+}, $d^3$: $S = \frac32$, $L = 2 + 1 + 0 = 3$,
\termsym{4}{F}{3/2}. \ce{Fe^2+}, $d^6$: five up, one down in $m_l = 2$:
$S = 2$, $L = 2$, more than half full, \termsym{5}{D}{4}. \ce{Ni^2+}, $d^8$:
$S = 1$, $L = 3$, \termsym{3}{F}{4}. Each agrees with the ground level that
atomic spectroscopy measures.
\end{example}

\begin{definition}[Spin--orbit coupling, fine structure]\label{def:b3:many-electron-atoms:spin-orbit}
\emph{Spin--orbit coupling}\index{spin--orbit coupling} is the interaction
between the magnetic moments of an electron's spin and of its orbital
motion, written $A\,\hat{\mathbf L}\cdot\hat{\mathbf S}/\hbar^2$ for a term. It
splits a term into its levels of different $J$: the \emph{fine
structure}\index{fine structure} of the term.
\end{definition}

\begin{proposition}[Landé interval rule]\label{prop:b3:many-electron-atoms:lande}
Within a term, $E(J) = E_0 + \frac{A}{2}[J(J+1) - L(L+1) - S(S+1)]$, so that
$E(J) - E(J - 1) = AJ$. $A > 0$ (normal) for a subshell less than half full,
$A < 0$ (inverted) for one more than half full.
\end{proposition}

\begin{proof}
$\hat{\mathbf J} = \hat{\mathbf L} + \hat{\mathbf S}$ gives $\hat{\mathbf J}^2 =
\hat{\mathbf L}^2 + \hat{\mathbf S}^2 + 2\hat{\mathbf L}\cdot\hat{\mathbf S}$, so
$\hat{\mathbf L}\cdot\hat{\mathbf S} = \frac12(\hat{\mathbf J}^2 - \hat{\mathbf L}^2 -
\hat{\mathbf S}^2)$, whose \omterm{def:b3:quantum-model-systems:operator}{eigenvalue} in a state of given $J$, $L$, $S$ is
$\frac{\hbar^2}{2}[J(J+1) - L(L+1) - S(S+1)]$. The difference between $J$ and
$J - 1$ is $\frac A2[J(J+1) - (J-1)J] = AJ$. The sign of $A$ is admitted: a
more-than-half-full subshell behaves as the corresponding number of positive
holes.
\end{proof}

\begin{example}[Testing the interval rule]\label{ex:b3:many-electron-atoms:lande-test}
% ledger: lev:C.3P1, lev:C.3P2, lev:O.3P1, lev:O.3P0
For ${}^3$P the rule predicts intervals in the ratio $J = 2 : 1$, that is 2.
Carbon: $(43.4 - 16.4)/16.4 = 1.65$. Oxygen ($2p^4$, inverted, \termsym{3}{P}{2}
lowest, then $J = 1$ at \qty{158.3}{cm^{-1}} and $J = 0$ at \qty{227.0}{cm^{-1}}):
$158.3/68.7 = 2.30$. The rule holds within 20\,\%: \omterm{def:b3:many-electron-atoms:russell-saunders}{Russell--Saunders coupling}
is a good description of light atoms, not an exact one.
\end{example}

\begin{omfigure}
\begin{tikzpicture}[font=\small]
  \draw[omstate] (0,0) -- (1.6,0) node[midway, below] {$3s\ {}^2$S$_{1/2}$};
  \draw[omstate] (0,3.0) -- (1.6,3.0);
  \node[anchor=east] at (-0.1,3.0) {$3p$};
  \draw[omstate] (2.8,2.75) -- (4.4,2.75) node[right] {${}^2$P$_{1/2}$};
  \draw[omstate] (2.8,3.25) -- (4.4,3.25) node[right] {${}^2$P$_{3/2}$};
  \draw[densely dotted, gray] (1.6,3.0) -- (2.8,2.75) (1.6,3.0) -- (2.8,3.25);
  \draw[omfluo] (3.3,2.73) -- (3.3,0.02) node[pos=0.55, left] {D$_1$};
  \draw[omfluo] (3.9,3.23) -- (3.9,0.02) node[pos=0.55, right] {D$_2$};
  \draw[omstate] (2.8,0) -- (4.4,0);
  \draw[densely dotted, gray] (1.6,0) -- (2.8,0);
  \draw[decorate, decoration={brace, amplitude=3pt}] (5.5,3.27) -- (5.5,2.73)
     node[midway, right=3pt, align=left] {\footnotesize \qty{17.2}{cm^{-1}}\\[-1pt]\footnotesize (not to scale)};
\end{tikzpicture}

{\small The sodium D lines. The $3p$ level is split by \omterm{def:b3:many-electron-atoms:spin-orbit}{spin--orbit coupling}
into ${}^2$P$_{1/2}$ and ${}^2$P$_{3/2}$; emission to the ground $3s$ level gives
two yellow lines, D$_1$ at \qty{589.76}{nm} and D$_2$ at \qty{589.16}{nm} (vacuum
wavelengths).}
\end{omfigure}

\begin{example}[The sodium doublet]\label{ex:b3:many-electron-atoms:sodium}
% ledger: lev:Na.3p1_2, lev:Na.3p3_2
The $3p$ levels of sodium lie at 16956.17 and \qty{16973.37}{cm^{-1}}: the D
lines are at $10^7/16956.17 = \qty{589.76}{nm}$ and \qty{589.16}{nm} in vacuum,
separated by \qty{17.20}{cm^{-1}}. For ${}^2$P, $E(\frac32) - E(\frac12) =
\frac32A$, so $A = \qty{11.47}{cm^{-1}}$.
\end{example}

\begin{proposition}[Selection rules for atoms]\label{prop:b3:many-electron-atoms:selection-atoms}
In \omterm{def:b3:many-electron-atoms:russell-saunders}{Russell--Saunders coupling}, an electric-dipole transition requires
$\Delta S = 0$, $\Delta L = 0, \pm1$, $\Delta J = 0, \pm1$ (but not $J = 0 \to J
= 0$), and a change of parity: for a one-electron jump, $\Delta l = \pm1$.
\end{proposition}
\begin{proof}[Status]
$\Delta S = 0$ is proved in \cref{ch:b3:electronic-spectroscopy}: the dipole
\omterm{def:b3:quantum-model-systems:operator}{operator} does not act on spin. The others follow from the vector character
of the dipole \omterm{def:b3:quantum-model-systems:operator}{operator} and are admitted here.
\end{proof}

\begin{omfigure}
\begin{tikzpicture}
\begin{axis}[width=0.62\linewidth, height=6.2cm, xmin=-0.5, xmax=3.5,
  ymin=158, ymax=173, axis y line=left, axis x line=none,
  ylabel={energy / $10^3$ cm$^{-1}$}, unbounded coords=jump, clip=false,
  tick label style={font=\small}, label style={font=\small}]
\addplot[omstate] table[x=x, y=E] {figdata/out/bachelor-3-atomic-levels-helium.dat};
\node[font=\small, anchor=west] at (axis cs:1.05,166.277) {$1s2s\ {}^1$S};
\node[font=\small, anchor=west] at (axis cs:1.05,171.135) {$1s2p\ {}^1$P};
\node[font=\small, anchor=west] at (axis cs:3.05,159.856) {$1s2s\ {}^3$S};
\node[font=\small, anchor=west] at (axis cs:3.05,169.087) {$1s2p\ {}^3$P};
\node[font=\small, anchor=south] at (axis cs:0.5,172.3) {singlets};
\node[font=\small, anchor=south] at (axis cs:2.5,172.3) {triplets};
\draw[omfluo] (axis cs:0.5,171.0) -- (axis cs:0.5,166.42) node[pos=0.5, left, font=\scriptsize] {\qty{2059}{nm}};
\draw[omfluo] (axis cs:2.5,168.95) -- (axis cs:2.5,160.0) node[pos=0.5, left, font=\scriptsize] {\qty{1083}{nm}};
\draw[densely dotted, gray] (axis cs:-0.3,159.856) -- (axis cs:2,159.856);
\draw[<->, omThm] (axis cs:-0.2,159.856) -- (axis cs:-0.2,166.277) node[pos=0.5, right, font=\scriptsize] {$2K$};
\end{axis}
\end{tikzpicture}

{\small The two ``heliums''. Singlet and triplet levels of the configurations
$1s2s$ and $1s2p$ at their measured energies above the ground state (the
${}^3$P term at the mean of its levels). Lines connect states of the same
family only; each triplet lies below its singlet, by $2K$.}
\end{omfigure}

\begin{history}[Two heliums and the exclusion principle]
Helium was named after the Sun in 1868 and found on Earth by William
Ramsay in 1895. Its two series of lines puzzled spectroscopists for thirty
years. In 1925 Wolfgang Pauli proposed that no two electrons share the
same quantum numbers, the same year as electron spin was proposed by George
Uhlenbeck and Samuel Goudsmit; in 1926 Werner Heisenberg showed that the
symmetry of the wavefunction alone splits helium into its singlet and
triplet families.
\end{history}

\begin{inthelab}[Resolving the sodium doublet]
A sodium lamp is placed before the slit of a grating spectrometer. With a
grating of 600 lines per millimetre, the first-order angles of the two lines
differ by about \num{0.02}\textdegree: a good spectrometer separates them, a
student prism does not. The lamp runs hot; it is handled only after it has
cooled.
\end{inthelab}

\section{Exercises}

\begin{exercise}[$\star$]\label{exo:b3:many-electron-atoms:1}
Write the \omterm{def:b3:many-electron-atoms:slater}{Slater determinant} of the ground state of lithium, $1s^22s$, with
the $2s$ electron of spin $\alpha$. Why is there no determinant for
$1s^3$?
\end{exercise}

\begin{exercise}[$\star$]\label{exo:b3:many-electron-atoms:2}
Count the determinants of the configurations $p^2$, $p^3$ and $d^2$. The terms
of $p^3$ are ${}^4$S, ${}^2$D, ${}^2$P and those of $d^2$ are ${}^3$F, ${}^3$P, ${}^1$G,
${}^1$D, ${}^1$S: check both counts.
\end{exercise}

\begin{exercise}[$\star$]\label{exo:b3:many-electron-atoms:3}
Give the ground term and level of N, O, F, \ce{Ti^2+} and \ce{Fe^3+} ($d^5$).
\end{exercise}

\begin{exercise}[$\star$]\label{exo:b3:many-electron-atoms:4}
List the levels of the terms ${}^2$D and ${}^3$P with their degeneracies $2J+1$,
and check that the degeneracies add up to $(2L+1)(2S+1)$.
\end{exercise}

\begin{exercise}[$\star\star$]\label{exo:b3:many-electron-atoms:5}
Carbon's ${}^3$P levels lie at 0, 16.4 and \qty{43.4}{cm^{-1}}. Compute the
ratio of the intervals and the spin--orbit constant $A$ from each interval.
What does the difference between the two values of $A$ say?
\end{exercise}

\begin{exercise}[$\star\star$]\label{exo:b3:many-electron-atoms:6}
The $3p$ levels of sodium lie at 16956.17 and \qty{16973.37}{cm^{-1}}. Compute
the vacuum wavelengths of the D lines, their separation in \unit{cm^{-1}} and
in meV, and $A$ for the $3p$ term.
\end{exercise}

\begin{exercise}[$\star\star$]\label{exo:b3:many-electron-atoms:7}
Which of these transitions of sodium are allowed in emission: $3p \to 3s$,
$3d \to 3s$, $3d \to 3p$, $4s \to 3s$, $4s \to 3p$? Justify each.
\end{exercise}

\begin{exercise}[$\star\star$]\label{exo:b3:many-electron-atoms:8}
% ledger: lev:He.2p3P2, lev:He.2p3P1, lev:He.2p3P0, lev:He.2p1P1
The $1s2p$ levels of helium are \termsym{3}{P}{2}, \termsym{3}{P}{1},
\termsym{3}{P}{0} at 169086.76, 169086.84, \qty{169087.83}{cm^{-1}} and
\termsym{1}{P}{1} at \qty{171134.89}{cm^{-1}}. Compute the mean energy of the
${}^3$P term, then $K(1s,2p)$ in eV. Compare with $K(1s,2s) = \qty{0.398}{eV}$
and explain the difference.
\end{exercise}

\begin{exercise}[$\star\star$]\label{exo:b3:many-electron-atoms:9}
Find the terms of $d^2$ by the method of this chapter, starting from the
largest $M_L$ (you may skip writing the whole table, but justify each term).
\end{exercise}

\begin{exercise}[$\star\star\star$]\label{exo:b3:many-electron-atoms:10}
Show that the triplet spatial function $\frac{1}{\sqrt2}[1s(1)2s(2) -
2s(1)1s(2)]$ vanishes when $\mathbf r_1 = \mathbf r_2$ and that the singlet
function does not. Relate this to the sign of $E_+ - E_-$.
\end{exercise}

\begin{exercise}[$\star\star\star$]\label{exo:b3:many-electron-atoms:11}
% ledger: lev:Ni2.3F3, lev:Ni2.3F2
For \ce{Ni^2+} ($d^8$) the ${}^3$F levels lie at 0 ($J = 4$), 1360.7 ($J = 3$)
and \qty{2269.6}{cm^{-1}} ($J = 2$). Explain the order, compute the ratio of
the intervals and compare with Landé's prediction. Give $A$ from the larger
interval.
\end{exercise}

\begin{exercise}[$\star\star\star$]\label{exo:b3:many-electron-atoms:12}
Show that the \omterm{def:b3:quantum-model-systems:hermitian}{expectation value} of $\hat{\mathbf L}\cdot\hat{\mathbf S}$, summed
over all the $(2L+1)(2S+1)$ states of a term, is zero, so that \omterm{def:b3:many-electron-atoms:spin-orbit}{spin--orbit
coupling} leaves the weighted mean of the levels unchanged. Check it on
carbon's ${}^3$P.
\end{exercise}

\section{Problem: Two Heliums}

\begin{problem}[{Weekend problem --- the singlet and triplet families of
helium: determinants, the terms of transition-metal ions, the selection
rules that separate the two families, and the exchange integral measured}]
\label{pb:b3:many-electron-atoms:1}
% ledger: lev:He.2s3S1, lev:He.2s1S0, lev:He.2p3P2, lev:He.2p3P1, lev:He.2p3P0, lev:He.2p1P1, lev:Ti2.3F3, lev:Ti2.3F4, lev:Ti2.3P0, lev:Ti2.3P1, lev:Ti2.3P2
Data (NIST, in \unit{cm^{-1}} above the ground state of each species). Helium:
$1s2s$ \termsym{3}{S}{1} 159855.97, \termsym{1}{S}{0} 166277.44; $1s2p$
\termsym{3}{P}{2,1,0} 169086.76, 169086.84, 169087.83, \termsym{1}{P}{1}
171134.89. \ce{Ti^2+}: \termsym{3}{F}{2,3,4} 0, 184.9, 420.4; \termsym{3}{P}{0,1,2}
10538.4, 10603.6, 10721.2. $\qty{1}{eV} = \qty{8065.54}{cm^{-1}}$.

\medskip\noindent\textbf{Part I --- \omterm{def:b3:many-electron-atoms:spin-orbital}{Spin-orbitals} and determinants.}
\begin{enumerate}
  \item List the four \omterm{def:b3:many-electron-atoms:spin-orbital}{spin-orbitals} available to the configuration $1s2s$.
  \item Write the four determinants with one electron in $1s$ and one in
        $2s$.
  \item Show that $|1s\alpha\,2s\alpha|$ is the product of an antisymmetric
        spatial function and the spin function $\alpha(1)\alpha(2)$.
  \item Combine $|1s\alpha\,2s\beta|$ and $|1s\beta\,2s\alpha|$ to obtain the
        $M_S = 0$ component of the same spatial function, and a fourth state.
  \item Which spatial function, symmetric or antisymmetric, goes with
        $S = 0$? With $S = 1$?
  \item Why are these called singlet and triplet?
\end{enumerate}

\medskip\noindent\textbf{Part II --- The terms of \ce{Ti^2+}.}
\begin{enumerate}[resume]
  \item Give the configuration of \ce{Ti^2+} and the number of its
        determinants.
  \item Find its ground term and level with Hund's rules; check against the
        data.
  \item Its other terms are ${}^3$P, ${}^1$G, ${}^1$D and ${}^1$S: check the count of
        states.
  \item Is the ${}^3$F \omterm{def:b3:many-electron-atoms:spin-orbit}{fine structure} normal or inverted? Test the Landé
        interval rule.
  \item Compute the weighted mean energies of the ${}^3$F and ${}^3$P terms.
  \item Their separation is $15B$, where $B$ is a Racah parameter of the
        ion (\cref{ch:b3:complex-spectra-magnetism}). Compute $B$.
\end{enumerate}

\medskip\noindent\textbf{Part III --- Why two families.}
\begin{enumerate}[resume]
  \item Which selection rule forbids a transition between a singlet and a
        triplet?
  \item The lowest triplet, $1s2s$ \termsym{3}{S}{1}, cannot emit to the ground
        state $1s^2$ \termsym{1}{S}{0}. Give two reasons. What is such a state
        called?
  \item Compute the wavelength of the line $1s2p\ {}^3$P $\to 1s2s\ {}^3$S (use
        the mean of the ${}^3$P levels).
  \item Compute the wavelength of $1s2p\ {}^1$P $\to 1s2s\ {}^1$S.
  \item Compute the wavelength of $1s2p\ {}^1$P $\to 1s^2\ {}^1$S. In which
        region of the spectrum is it?
  \item Explain why nineteenth-century spectroscopists, seeing these lines,
        believed in two different gases.
\end{enumerate}

\medskip\noindent\textbf{Part IV --- The \omterm{def:b3:many-electron-atoms:exchange}{exchange integral}, measured.}
\begin{enumerate}[resume]
  \item Write the first-order energies of the $1s2s$ singlet and triplet in
        terms of $E_{1s}$, $E_{2s}$, $J$ and $K$.
  \item Which is lower, and why, in terms of the electrons' positions?
  \item Compute the $1s2s$ singlet--triplet gap in \unit{cm^{-1}} and eV.
  \item Compute $K(1s,2p)$ in eV from the $1s2p$ levels.
  \item Why is $K(1s,2s)$ larger than $K(1s,2p)$?
  \item Is Hund's first rule borne out by helium's excited states?
  \item State the result: the \omterm{def:b3:many-electron-atoms:exchange}{exchange integral} $K(1s,2s)$ of helium, in
        eV.
\end{enumerate}
\end{problem}
